\chapter{Gaussian concentration and the Borell--TIS inequality}
\label{chap:pre_concentration}

This chapter proves that the supremum $M = \sup_{x\in\Xs}\ip{x}{G}$ of a linear Gaussian
process over a compact index set concentrates around its mean at the scale
$\sigma := \sup_{x\in\Xs}\sqrt{x^\top S x}$ of the largest individual standard deviation
(the Borell--TIS inequality, Theorem~\ref{thm:borell_tis}), in the form used by the
fixed-design upper bound (Corollaries~\ref{cor:sup_bound_whp}
and~\ref{cor:sup_bound_symmetric}). The route is the Maurey--Pisier argument
\cite{pisier1986probabilistic} (see also \cite[Section~5.4]{vershynin2018high} and
\cite[Chapter~5]{boucheron2013concentration}): a $C^1$ function $f$ of a standard Gaussian
vector with $\norm{\nabla f} \le L$ has a sub-Gaussian moment generating function with
constant $\cG L^2$, where
\[
  \cG := \pi^2/4 ,
\]
and the supremum of a finite family of linear forms is approximated by the smooth
``log-sum-exp'' function. The optimal constant is $1$ (Gaussian isoperimetry, or the
Gaussian log-Sobolev inequality, see \cite{adler2007random}); the factor $\cG$ is the price
of an argument that uses only Jensen's inequality, the fundamental theorem of calculus, Fubini
and the rotation invariance of Gaussian pairs, and it propagates as an explicit absolute
constant into the upper bounds of Part~\ref{part:main}.

\textbf{What Mathlib provides.} The structure \texttt{ProbabilityTheory.HasSubgaussianMGF X c μ}
(\texttt{Mathlib/Probability/Moments/SubGaussian.lean}): integrability of $e^{tX}$ for all
$t$ together with $\E e^{tX} \le e^{ct^2/2}$; its Chernoff bound
\texttt{HasSubgaussianMGF.measure\_ge\_le}, and \texttt{measure\_le\_le} in LML
(\texttt{ForMathlib/Probability/Moments/SubGaussian.lean}). Jensen's inequality
\texttt{ConvexOn.map\_integral\_le} and \texttt{ConvexOn.map\_average\_le}
(\texttt{Mathlib/Analysis/Convex/Integral.lean}) with \texttt{convexOn\_exp}. The fundamental
theorem of calculus \texttt{intervalIntegral.integral\_eq\_sub\_of\_hasDerivAt}, the chain
rule \texttt{HasFDerivAt.comp\_hasDerivAt}, the mean value inequality
\texttt{lipschitzWith\_of\_nnnorm\_fderiv\_le}, Fubini \texttt{MeasureTheory.integral\_prod}
and \texttt{integral\_integral\_swap}, dominated convergence
\texttt{tendsto\_integral\_of\_dominated\_convergence}, Fernique's theorem
\texttt{IsGaussian.exists\_integrable\_exp\_sq}, the rotation invariance of Gaussian pairs
(Lemma~\ref{lem:gauss_rotation}, Mathlib \texttt{stdGaussian\_map} and
\texttt{IsGaussian.map\_rotation\_eq\_self}), and the Gaussian MGF \texttt{mgf\_gaussianReal}.

\textbf{What is missing.} Everything specific to concentration: the Maurey--Pisier inequality
(Lemma~\ref{lem:maurey_pisier}), the log-sum-exp smoothing (Lemma~\ref{lem:logsumexp_props}),
the passage from finite to compact index sets (Lemma~\ref{lem:compact_sup_subgaussian}) and
the Borell--TIS inequality itself. Mathlib has no Gaussian concentration inequality for
Lipschitz functions at the time of writing.

\textbf{Lean remark.} Sub-Gaussianity of $M - \E M$ is stated with Mathlib's
\texttt{HasSubgaussianMGF (fun ω ↦ M ω - P[M]) (cG * σ²) P}, whose constant is an
\texttt{ℝ≥0}; all constants below are therefore nonnegative reals coerced to \texttt{ℝ≥0}.
Gradients are \texttt{gradient f x} in \texttt{EuclideanSpace ℝ (Fin d)}, with
$\ip{\nabla f(x)}{v} = f'(x)(v)$ (\texttt{HasGradientAt} / \texttt{HasFDerivAt}), and
``$\norm{\nabla f} \le L$ everywhere'' is \texttt{∀ x, ‖gradient f x‖ ≤ L} together with
differentiability; the class $C^1$ is \texttt{ContDiff ℝ 1 f}.

\section{Preliminaries}
\label{sec:pre_concentration_prelim}

\begin{definition}[The Gaussian concentration constant]\label{def:pc_gaussian_constant}
$\cG := \pi^2/4$ ($\approx 2.4674$).
\end{definition}

\begin{lemma}[MGF of a linear form of a standard Gaussian vector]\label{lem:gauss_mgf_inner}
\uses{def:pg_gaussian_vector}
Let $g' \sim \Normal(0, I_d)$ and $a \in \R^d$. Then $e^{\ip{a}{g'}}$ is integrable and
$\E e^{\ip{a}{g'}} = e^{\norm{a}^2/2}$. More generally, for $s \in \R$,
$\E e^{s\ip{a}{g'}} = e^{s^2\norm{a}^2/2}$.
\end{lemma}

\begin{proof}
\uses{lem:gauss_linear_image, lem:gauss_1d_moments}
$\ip{a}{g'} \sim \Normal(0,\norm{a}^2)$ by Lemma~\ref{lem:gauss_linear_image}, and
Lemma~\ref{lem:gauss_1d_moments} (Mathlib \texttt{mgf\_gaussianReal} at the point $s$, with
\texttt{integrable\_exp\_mul\_gaussianReal}) gives the value. Alternatively, by
Lemma~\ref{lem:gauss_pi}, $\ip{a}{g'} = \sum_ia_ig'_i$ with independent coordinates, and
$\E\prod_ie^{sa_ig'_i} = \prod_ie^{s^2a_i^2/2}$ (\texttt{iIndepFun.mgf\_sum}).
\end{proof}

\begin{lemma}[Exponential moments of Gaussian norms]\label{lem:pc_exp_norm_integrable}
\uses{def:pg_gaussian_vector}
Let $G \sim \Normal(\mu, S)$ in $\R^n$ and $c \ge 0$. Then $e^{c\norm{G}}$ is integrable.
Consequently, if $F : \R^n \to \R$ is measurable with $\abs{F(x)} \le A + B\norm{x}$ for
constants $A, B \ge 0$, then $F(G)$ is integrable and $e^{\lambda F(G)}$ is integrable for every
$\lambda \in \R$.
\end{lemma}

\begin{proof}
\uses{lem:gauss_1d_moments, lem:gauss_pi}
By Fernique's theorem (Mathlib \texttt{IsGaussian.exists\_integrable\_exp\_sq}) there is
$C > 0$ with $e^{C\norm{x}^2}$ integrable under the law of $G$; since
$c\norm{x} \le C\norm{x}^2 + c^2/(4C)$ (AM--GM), $e^{c\norm{G}} \le e^{c^2/(4C)}e^{C\norm{G}^2}$
is integrable. (Elementary alternative for $G = \mu + S^{1/2}g$: $\norm{G} \le \norm{\mu} +
\norm{S^{1/2}}_{\mathrm{op}}\sum_i\abs{g_i}$, $e^{c'\abs{g_i}} \le e^{c'g_i} + e^{-c'g_i}$, and
by independence of the coordinates (Lemma~\ref{lem:gauss_pi}) and
Lemma~\ref{lem:gauss_1d_moments} the expectation of the product is a finite product of finite
numbers.) For $F$: $\abs{F(G)} \le A + B\norm{G}$ is integrable, and
$e^{\lambda F(G)} \le e^{\abs\lambda A}e^{\abs\lambda B\norm{G}}$ is integrable by the first
part with $c = \abs{\lambda}B$.
\end{proof}

\begin{lemma}[Bounded gradient implies Lipschitz and linear growth]\label{lem:pc_lipschitz_of_grad}
Let $f : \R^d \to \R$ be differentiable with $\norm{\nabla f(x)} \le L$ for every $x$, where
$L \ge 0$. Then $\abs{f(x) - f(y)} \le L\norm{x-y}$ for all $x,y$, $\abs{f(x)} \le \abs{f(0)} +
L\norm{x}$, and $\abs{\ip{\nabla f(x)}{v}} \le L\norm{v}$ for all $x, v$. In particular, for
$g \sim \Normal(0,I_d)$, $f(g)$ is integrable and $e^{\lambda f(g)}$ is integrable for every
$\lambda\in\R$.
\end{lemma}

\begin{proof}
\uses{lem:pc_exp_norm_integrable}
The Lipschitz bound is the mean value inequality (Mathlib
\texttt{lipschitzWith\_of\_nnnorm\_fderiv\_le}, the operator norm of $f'(x)$ being
$\norm{\nabla f(x)}$ in an inner product space); $y = 0$ gives the growth bound; the last
inequality is Cauchy--Schwarz. Integrability is Lemma~\ref{lem:pc_exp_norm_integrable} with
$A = \abs{f(0)}$, $B = L$.
\end{proof}

\begin{lemma}[The log-sum-exp smoothing of a finite maximum]\label{lem:logsumexp_props}
Let $m \ge 1$, $x_1,\dots,x_m \in \R^d$, $\beta > 0$ and
\[
  f_\beta(v) := \frac1\beta\log\sum_{i\le m}e^{\beta\ip{x_i}{v}}, \qquad v \in \R^d .
\]
\begin{enumerate}
  \item[(i)] $f_\beta$ is $C^\infty$ on $\R^d$, with gradient
        $\nabla f_\beta(v) = \sum_{i\le m}p_i(v)\,x_i$ where
        $p_i(v) := e^{\beta\ip{x_i}{v}}/\sum_{j\le m}e^{\beta\ip{x_j}{v}}$ satisfies $p_i(v) \ge 0$,
        $\sum_ip_i(v) = 1$; hence $\norm{\nabla f_\beta(v)} \le \max_{i\le m}\norm{x_i}$ for all $v$.
  \item[(ii)] For every $v$,
        $\max_{i\le m}\ip{x_i}{v} \le f_\beta(v) \le \max_{i\le m}\ip{x_i}{v} + \frac{\log m}{\beta}$.
\end{enumerate}
\end{lemma}

\begin{proof}
(i) $v \mapsto \sum_ie^{\beta\ip{x_i}{v}}$ is a finite sum of compositions of $\exp$ with
linear forms, hence $C^\infty$ and strictly positive; $\log$ is $C^\infty$ on $(0,\infty)$
(Mathlib \texttt{ContDiff.log}, \texttt{contDiff\_exp}, \texttt{ContDiff.inner}). The chain
rule gives $\nabla f_\beta(v) = \frac1\beta\cdot\frac{\sum_i\beta e^{\beta\ip{x_i}{v}}x_i}
{\sum_je^{\beta\ip{x_j}{v}}} = \sum_ip_i(v)x_i$. The $p_i(v)$ are positive with sum $1$, so by
the triangle inequality $\norm{\nabla f_\beta(v)} \le \sum_ip_i(v)\norm{x_i} \le \max_i\norm{x_i}$.
(ii) Let $M(v) := \max_i\ip{x_i}{v}$. Since $\exp$ and $\log$ are increasing,
$e^{\beta M(v)} \le \sum_ie^{\beta\ip{x_i}{v}} \le m\,e^{\beta M(v)}$; take logarithms and
divide by $\beta$.
\end{proof}

\section{The Maurey--Pisier argument}
\label{sec:pre_concentration_mp}

Throughout this section $f : \R^d \to \R$ is of class $C^1$ with $\norm{\nabla f(x)} \le L$
for all $x$, and $g, g'$ are independent $\Normal(0,I_d)$ vectors on $(\Omega,\Prob)$. For
$\theta \in \R$ we write, as in Lemma~\ref{lem:gauss_rotation},
\[
  g_\theta := g\sin\theta + g'\cos\theta, \qquad g'_\theta := g\cos\theta - g'\sin\theta,
\]
so that $g_0 = g'$, $g_{\pi/2} = g$ and $\frac{d}{d\theta}g_\theta = g'_\theta$.

\begin{lemma}[Interpolation along a quarter circle]\label{lem:pc_interpolation_ftc}
Let $f : \R^d \to \R$ be of class $C^1$ and $u, v \in \R^d$. Set
$\gamma(\theta) := u\sin\theta + v\cos\theta$ and $\gamma'(\theta) := u\cos\theta - v\sin\theta$.
Then $\theta \mapsto f(\gamma(\theta))$ is differentiable on $\R$ with derivative
$\theta \mapsto \ip{\nabla f(\gamma(\theta))}{\gamma'(\theta)}$, this derivative is continuous,
and
\[
  f(u) - f(v) = \int_0^{\pi/2}\ip{\nabla f(\gamma(\theta))}{\gamma'(\theta)}\,d\theta .
\]
\end{lemma}

\begin{proof}
$\gamma$ is differentiable with $\gamma'(\theta) = u\cos\theta - v\sin\theta$
(\texttt{HasDerivAt.smul\_const} for $\sin$ and $\cos$, \texttt{HasDerivAt.add}). The chain
rule (\texttt{HasFDerivAt.comp\_hasDerivAt}) gives that $f\circ\gamma$ has derivative
$f'(\gamma(\theta))(\gamma'(\theta)) = \ip{\nabla f(\gamma(\theta))}{\gamma'(\theta)}$ at
$\theta$. This derivative is continuous in $\theta$ as a composition of continuous maps
($\nabla f$ is continuous since $f$ is $C^1$), hence interval integrable on $[0,\pi/2]$
(\texttt{ContinuousOn.intervalIntegrable}). The fundamental theorem of calculus
(\texttt{intervalIntegral.integral\_eq\_sub\_of\_hasDerivAt}) gives
$\int_0^{\pi/2}(f\circ\gamma)' = f(\gamma(\pi/2)) - f(\gamma(0)) = f(u) - f(v)$.
\end{proof}

\begin{lemma}[{Jensen's inequality for the uniform measure on $[0,\pi/2]$}]
\label{lem:pc_jensen_exp_average}
Let $h : [0,\pi/2] \to \R$ be continuous and $s \in \R$. Then
\[
  \exp\Big(s\int_0^{\pi/2}h(\theta)\,d\theta\Big)
  \le \frac2\pi\int_0^{\pi/2}\exp\Big(\frac{\pi s}{2}\,h(\theta)\Big)d\theta .
\]
\end{lemma}

\begin{proof}
Let $\nu := \frac2\pi\,\mathrm{Leb}|_{[0,\pi/2]}$, a probability measure, and
$k(\theta) := \frac{\pi s}{2}h(\theta)$, which is continuous, hence bounded on the compact
interval; so $k$ and $e^{k}$ are $\nu$-integrable. Jensen's inequality for the convex
continuous function $\exp$ on the closed convex set $\R$ (Mathlib
\texttt{ConvexOn.map\_integral\_le} with \texttt{convexOn\_exp}, \texttt{continuous\_exp},
\texttt{isClosed\_univ}) gives $\exp(\int k\,d\nu) \le \int e^{k}\,d\nu$. Finally
$\int k\,d\nu = \frac2\pi\cdot\frac{\pi s}{2}\int_0^{\pi/2}h = s\int_0^{\pi/2}h$ and
$\int e^k\,d\nu = \frac2\pi\int_0^{\pi/2}e^{k(\theta)}d\theta$.
\end{proof}

\begin{lemma}[MGF of the rotated gradient term]\label{lem:pc_rotated_mgf_bound}
\uses{def:pg_gaussian_vector, lem:pc_lipschitz_of_grad}
Let $f$ be $C^1$ with $\norm{\nabla f} \le L$ everywhere, $g, g'$ independent
$\Normal(0,I_d)$ vectors, $s \in \R$ and $\theta \in \R$. Then
$\omega \mapsto \exp(s\ip{\nabla f(g_\theta)}{g'_\theta})$ is integrable and
\[
  \E\exp\big(s\ip{\nabla f(g_\theta)}{g'_\theta}\big)
  = \E\exp\big(s\ip{\nabla f(g)}{g'}\big)
  = \E\exp\Big(\frac{s^2\norm{\nabla f(g)}^2}{2}\Big) \le \exp\Big(\frac{s^2L^2}{2}\Big).
\]
\end{lemma}

\begin{proof}
\uses{lem:gauss_rotation, lem:gauss_mgf_inner, lem:pc_lipschitz_of_grad, lem:pc_exp_norm_integrable}
Let $\Phi(u,v) := \exp(s\ip{\nabla f(u)}{v})$, a continuous function on $\R^d\times\R^d$ with
$0 < \Phi(u,v) \le \exp(\abs{s}L\norm{v})$ (Lemma~\ref{lem:pc_lipschitz_of_grad}). By
Lemma~\ref{lem:gauss_rotation}(ii), $(g_\theta, g'_\theta)$ has the same law as $(g,g')$, so
$\E\Phi(g_\theta,g'_\theta) = \E\Phi(g,g')$ (equality of laws; both sides are finite by the
domination and Lemma~\ref{lem:pc_exp_norm_integrable}). The law of $(g,g')$ is the product
$\Normal(0,I_d)\otimes\Normal(0,I_d)$ (independence), and $\Phi$ is integrable for it (dominated
by $\exp(\abs sL\norm{v})$, which is integrable in $v$ and constant in $u$), so Fubini
(\texttt{MeasureTheory.integral\_prod}) gives
\[
  \E\Phi(g,g') = \int\Big(\int\exp\big(\ip{s\nabla f(u)}{v}\big)\,\Normal(0,I_d)(dv)\Big)\Normal(0,I_d)(du)
  = \int\exp\Big(\frac{s^2\norm{\nabla f(u)}^2}{2}\Big)\Normal(0,I_d)(du),
\]
by Lemma~\ref{lem:gauss_mgf_inner} with $a = s\nabla f(u)$. The last integrand is at most
$\exp(s^2L^2/2)$ pointwise, and $\Normal(0,I_d)$ is a probability measure.
\end{proof}

\begin{lemma}[Maurey--Pisier Gaussian concentration]\label{lem:maurey_pisier}
\uses{def:pg_gaussian_vector, def:pc_gaussian_constant}
Let $L \ge 0$ and let $f : \R^d \to \R$ be of class $C^1$ with $\norm{\nabla f(x)} \le L$ for
every $x \in \R^d$. Let $g \sim \Normal(0, I_d)$. Then $f(g)$ is integrable, and for every
$\lambda \in \R$ the variable $e^{\lambda(f(g) - \E f(g))}$ is integrable with
\[
  \E\exp\big(\lambda(f(g) - \E f(g))\big) \le \exp\Big(\frac{\pi^2\lambda^2L^2}{8}\Big)
  = \exp\Big(\frac{\cG L^2\,\lambda^2}{2}\Big).
\]
In Mathlib's terms, $f(g) - \E f(g)$ satisfies \texttt{HasSubgaussianMGF} with constant
$\cG L^2$.
\end{lemma}

\begin{proof}
\uses{lem:pc_lipschitz_of_grad, lem:pc_exp_norm_integrable, lem:pc_interpolation_ftc, lem:pc_jensen_exp_average, lem:pc_rotated_mgf_bound, lem:gauss_rotation}
\emph{Step 0 (integrability).} By Lemma~\ref{lem:pc_lipschitz_of_grad}, $f$ is $L$-Lipschitz,
$\abs{f(x)} \le \abs{f(0)} + L\norm{x}$, $f(g)$ is integrable, and $e^{\lambda f(g)}$ is
integrable for every $\lambda$; hence so is $e^{\lambda(f(g)-\E f(g))}$. Fix $\lambda \in \R$
(for $\lambda = 0$ there is nothing to prove). Enlarging the probability space if necessary
(all quantities depend only on the law of $g$), let $g'$ be an independent copy of $g$; the
law of $(g,g')$ is $\gamma\otimes\gamma$ with $\gamma := \Normal(0,I_d)$.

\emph{Step 1 (symmetrization by Jensen in $g'$).} For fixed $u \in \R^d$, the variable
$\lambda(f(u) - f(g'))$ is integrable, and $e^{\lambda(f(u)-f(g'))} \le
e^{\abs\lambda(\abs{f(u)} + \abs{f(0)} + L\norm{g'})}$ is integrable
(Lemma~\ref{lem:pc_exp_norm_integrable}); Jensen's inequality for $\exp$
(\texttt{ConvexOn.map\_integral\_le}, \texttt{convexOn\_exp}) gives
\[
  \exp\big(\lambda(f(u) - \E f(g'))\big) \le \E\exp\big(\lambda(f(u) - f(g'))\big).
\]
Since $\E f(g') = \E f(g)$, evaluating at $u = g$ and integrating in $g$ (Fubini for
$\gamma\otimes\gamma$, \texttt{MeasureTheory.integral\_prod}; the function
$(u,v) \mapsto e^{\lambda(f(u)-f(v))}$ is dominated by
$e^{2\abs\lambda\abs{f(0)}}e^{\abs\lambda L\norm{u}}e^{\abs\lambda L\norm{v}}$, integrable on the
product by Lemma~\ref{lem:pc_exp_norm_integrable}),
\[
  \E\exp\big(\lambda(f(g) - \E f(g))\big) \le \E\exp\big(\lambda(f(g) - f(g'))\big).
\]

\emph{Step 2 (interpolation).} Pointwise in $\omega$, Lemma~\ref{lem:pc_interpolation_ftc}
with $u = g(\omega)$, $v = g'(\omega)$ gives
\[
  f(g) - f(g') = \int_0^{\pi/2}h_\omega(\theta)\,d\theta, \qquad
  h_\omega(\theta) := \ip{\nabla f(g_\theta)}{g'_\theta},
\]
where $\theta \mapsto h_\omega(\theta)$ is continuous.

\emph{Step 3 (Jensen for the uniform measure on $[0,\pi/2]$).} By
Lemma~\ref{lem:pc_jensen_exp_average} with $s = \lambda$, pointwise in $\omega$,
\[
  \exp\big(\lambda(f(g) - f(g'))\big)
  \le \frac2\pi\int_0^{\pi/2}\exp\Big(\frac{\pi\lambda}{2}\,h_\omega(\theta)\Big)d\theta .
\]

\emph{Step 4 (Fubini and rotation invariance).} The function
$(\omega,\theta) \mapsto \exp(\frac{\pi\lambda}{2}h_\omega(\theta))$ is measurable (a continuous
function of $(g(\omega), g'(\omega), \theta)$) and satisfies, by
Lemma~\ref{lem:pc_lipschitz_of_grad} and $\norm{g'_\theta} \le \norm{g} + \norm{g'}$,
\[
  0 < \exp\Big(\frac{\pi\lambda}{2}h_\omega(\theta)\Big)
  \le \exp\Big(\frac{\pi\abs\lambda L}{2}\big(\norm{g(\omega)} + \norm{g'(\omega)}\big)\Big) =: D(\omega),
\]
where $D$ does not depend on $\theta$ and is integrable on $\Omega$ (product of two
independent integrable variables, Lemma~\ref{lem:pc_exp_norm_integrable}); hence the function
is integrable on $\Omega\times[0,\pi/2]$ and Fubini (\texttt{integral\_integral\_swap})
applies:
\[
  \E\exp\big(\lambda(f(g) - f(g'))\big)
  \le \frac2\pi\int_0^{\pi/2}\E\exp\Big(\frac{\pi\lambda}{2}\ip{\nabla f(g_\theta)}{g'_\theta}\Big)d\theta .
\]
For each fixed $\theta$, Lemma~\ref{lem:pc_rotated_mgf_bound} with $s = \pi\lambda/2$ gives
\[
  \E\exp\Big(\frac{\pi\lambda}{2}\ip{\nabla f(g_\theta)}{g'_\theta}\Big)
  \le \exp\Big(\frac{(\pi\lambda/2)^2L^2}{2}\Big) = \exp\Big(\frac{\pi^2\lambda^2L^2}{8}\Big).
\]

\emph{Step 5 (conclusion).} Integrating the constant bound over $\theta \in [0,\pi/2]$ and
multiplying by $2/\pi$,
\[
  \E\exp\big(\lambda(f(g) - \E f(g))\big) \le \frac2\pi\cdot\frac\pi2\cdot\exp\Big(\frac{\pi^2\lambda^2L^2}{8}\Big)
  = \exp\Big(\frac{\pi^2\lambda^2L^2}{8}\Big).
\]
Finally $\frac{\pi^2L^2}{8}\lambda^2 = \frac{(\pi^2/4)L^2\,\lambda^2}{2} = \frac{\cG L^2\lambda^2}{2}$,
which together with Step~0 is the statement \texttt{HasSubgaussianMGF} with constant
$\cG L^2$.
\end{proof}

\begin{remark}[On the constant]\label{rem:pc_constant}
The factor $\pi^2/8$ in the exponent comes from Step~3: Jensen's inequality replaces
$\lambda\int_0^{\pi/2}h$ by an average of $\exp(\frac\pi2\lambda h(\theta))$, and the length
$\pi/2$ of the interval is squared in the Gaussian MGF of Step~4. With the standard
Gaussian interpolation $g_t = \sqrt t\,g + \sqrt{1-t}\,g'$ one would need the derivative
$\frac{d}{dt}g_t$, which is singular at the endpoints; the quarter-circle parametrization
avoids this at the cost of the constant. Sharper routes (Gaussian log-Sobolev + Herbst, or
the Gaussian isoperimetric inequality) give $\cG = 1$ \cite{boucheron2013concentration,
adler2007random}, and the Borell--TIS inequality with $\cG = 1$ is the form usually quoted;
none of the results of Part~\ref{part:main} is sensitive to the value of $\cG$ beyond the
absolute constants.
\end{remark}

\section{Suprema of linear Gaussian processes}
\label{sec:pre_concentration_sup}

\begin{lemma}[Concentration of a finite maximum of Gaussian linear forms]
\label{lem:finite_sup_subgaussian}
\uses{def:pg_gaussian_vector, def:pc_gaussian_constant}
Let $m \ge 1$, $x_1,\dots,x_m \in \R^d$, $S \in \PSD$ and $G \sim \Normal(0,S)$. Let
$M := \max_{i\le m}\ip{x_i}{G}$ and $\sigma^2 := \max_{i\le m}x_i^\top Sx_i$. Then $M$ is
integrable and $M - \E M$ has a sub-Gaussian MGF with constant $\cG\sigma^2$: for every
$\lambda\in\R$, $e^{\lambda(M - \E M)}$ is integrable and
$\E e^{\lambda(M-\E M)} \le \exp(\cG\sigma^2\lambda^2/2)$.
\end{lemma}

\begin{proof}
\uses{lem:gauss_law_of_cov, lem:maurey_pisier, lem:logsumexp_props, lem:pc_exp_norm_integrable}
The statement only involves the law of $M$, a measurable function of $G$, so by
Lemma~\ref{lem:gauss_law_of_cov}(iii) we may assume $G = S^{1/2}g$ with $g \sim \Normal(0,I_d)$
(Mathlib \texttt{HasSubgaussianMGF.congr\_identDistrib} / \texttt{of\_map} transport the
conclusion). Since $S^{1/2}$ is symmetric, $\ip{x_i}{S^{1/2}g} = \ip{y_i}{g}$ with
$y_i := S^{1/2}x_i$, and $\norm{y_i}^2 = x_i^\top Sx_i \le \sigma^2$. Thus $M = F(g)$ with
$F(v) := \max_{i\le m}\ip{y_i}{v}$, and $\abs{F(v)} \le \sigma\norm{v}$, so $M$ and
$e^{\lambda M}$ are integrable (Lemma~\ref{lem:pc_exp_norm_integrable}).

Fix $\lambda \in \R$ and $\beta > 0$, and let $f_\beta$ be the log-sum-exp function of
Lemma~\ref{lem:logsumexp_props} for the vectors $y_1,\dots,y_m$. It is $C^\infty$ with
$\norm{\nabla f_\beta} \le \max_i\norm{y_i} \le \sigma$, so Lemma~\ref{lem:maurey_pisier} with
$L = \sigma$ gives $\E e^{\lambda(f_\beta(g) - \E f_\beta(g))} \le \exp(\cG\sigma^2\lambda^2/2)$.
By Lemma~\ref{lem:logsumexp_props}(ii), $0 \le f_\beta(g) - M \le (\log m)/\beta$ pointwise,
hence also $0 \le \E f_\beta(g) - \E M \le (\log m)/\beta$, and therefore
\[
  \abs{(M - \E M) - (f_\beta(g) - \E f_\beta(g))} \le \frac{2\log m}{\beta}
  \quad\text{pointwise, so}\quad
  \E e^{\lambda(M - \E M)} \le e^{2\abs\lambda(\log m)/\beta}\,\E e^{\lambda(f_\beta(g) - \E f_\beta(g))}
  \le e^{2\abs\lambda(\log m)/\beta}\exp\Big(\frac{\cG\sigma^2\lambda^2}{2}\Big).
\]
The left-hand side does not depend on $\beta$; letting $\beta \to \infty$ (or taking the
infimum over $\beta > 0$: a real number bounded by $a e^{\eps}$ for every $\eps > 0$ is bounded
by $a$) gives $\E e^{\lambda(M-\E M)} \le \exp(\cG\sigma^2\lambda^2/2)$. No limit theorem for
integrals is needed here.
\end{proof}

\begin{lemma}[Concentration of the supremum over a compact index set]
\label{lem:compact_sup_subgaussian}
\uses{def:pg_gaussian_vector, def:pc_gaussian_constant, lem:sup_countable_dense}
Let $\Xs \subset \R^d$ be nonempty and compact, $S \in \PSD$ and $G \sim \Normal(0,S)$. Let
$M := \sup_{x\in\Xs}\ip{x}{G}$ (a random variable by Lemma~\ref{lem:sup_countable_dense})
and $\sigma^2 := \sup_{x\in\Xs}x^\top Sx$ (finite, attained). Then $M$ is integrable and
$M - \E M$ has a sub-Gaussian MGF with constant $\cG\sigma^2$: for every $\lambda\in\R$,
$e^{\lambda(M-\E M)}$ is integrable and $\E e^{\lambda(M - \E M)} \le \exp(\cG\sigma^2\lambda^2/2)$.
\end{lemma}

\begin{proof}
\uses{lem:sup_countable_dense, lem:finite_sup_subgaussian, lem:pc_exp_norm_integrable}
Let $R := \max_{x\in\Xs}\norm{x}$ and let $(q_n)_{n\in\N}$ enumerate a countable dense subset
of $\Xs$ (Lemma~\ref{lem:sup_countable_dense}(i)). $\sigma^2$ is the maximum of the continuous
function $x\mapsto x^\top Sx = \norm{S^{1/2}x}^2$ on the compact $\Xs$. Define the partial
maxima $M_n := \max_{i\le n}\ip{q_i}{G}$, $n \in \N$. By Lemma~\ref{lem:sup_countable_dense}(ii)
applied pointwise to $v = G(\omega)$, $M_n \uparrow M$ everywhere, and
$\abs{M_n} \le R\norm{G}$, $\abs{M} \le R\norm{G}$.

\emph{Bound for each $n$.} Lemma~\ref{lem:finite_sup_subgaussian} applied to $q_0,\dots,q_n$
gives, with $\sigma_n^2 := \max_{i\le n}q_i^\top Sq_i \le \sigma^2$, that
$\E e^{\lambda(M_n - \E M_n)} \le \exp(\cG\sigma_n^2\lambda^2/2) \le \exp(\cG\sigma^2\lambda^2/2)$
for every $\lambda$.

\emph{Passage to the limit.} Fix $\lambda\in\R$. The dominating function
$e^{\abs\lambda R\norm{G}}$ is integrable (Lemma~\ref{lem:pc_exp_norm_integrable}) and
$e^{\lambda M_n} \le e^{\abs\lambda R\norm{G}}$ for all $n$, while $e^{\lambda M_n} \to e^{\lambda M}$
pointwise (continuity of $\exp$). Dominated convergence
(\texttt{tendsto\_integral\_of\_dominated\_convergence}) gives
$\E e^{\lambda M_n} \to \E e^{\lambda M}$, and $e^{\lambda M}$ is integrable. Likewise
$\abs{M_n} \le R\norm{G}$, $M_n \to M$, and $R\norm{G}$ is integrable, so
$\E M_n \to \E M$ (dominated convergence again, or monotone convergence
\texttt{integral\_tendsto\_of\_tendsto\_of\_monotone} since $M_n$ is nondecreasing); in
particular $M$ is integrable. Therefore
\[
  \E e^{\lambda(M - \E M)} = e^{-\lambda\E M}\,\E e^{\lambda M}
  = \lim_{n\to\infty}e^{-\lambda\E M_n}\,\E e^{\lambda M_n}
  = \lim_{n\to\infty}\E e^{\lambda(M_n - \E M_n)} \le \exp\Big(\frac{\cG\sigma^2\lambda^2}{2}\Big),
\]
as a limit of numbers bounded by the right-hand side (\texttt{le\_of\_tendsto}). Integrability
of $e^{\lambda(M-\E M)} = e^{-\lambda\E M}e^{\lambda M}$ was shown above.
\end{proof}

\textbf{Lean remark.} The supremum over $\Xs$ is encoded as an \texttt{iSup} over the dense
sequence \texttt{q : ℕ → Xs} (Lemma~\ref{lem:sup_countable_dense}), so that $M_n$ is
\texttt{⨆ i : Fin (n+1), ⟪q i, G ω⟫} (or a \texttt{Finset.sup'} over \texttt{Finset.range}),
and the two limits are \texttt{Filter.Tendsto} statements along \texttt{atTop}. Lemma
\ref{lem:finite_sup_subgaussian} is stated for an indexed family \texttt{x : Fin m → E}, $m \ge 1$,
and the constant $\sigma^2$ is \texttt{⨆ i, ⟪x i, S (x i)⟫} coerced to \texttt{ℝ≥0}.

\begin{theorem}[Borell--TIS inequality for linear Gaussian processes]\label{thm:borell_tis}
\uses{def:pg_gaussian_vector, def:pc_gaussian_constant, lem:sup_countable_dense}
Let $\Xs \subset \R^d$ be nonempty and compact, $S \in \PSD$, $G \sim \Normal(0,S)$,
$M := \sup_{x\in\Xs}\ip{x}{G}$ and $\sigma^2 := \sup_{x\in\Xs}x^\top Sx$, assumed positive.
Then $M$ is integrable and for every $u \ge 0$,
\[
  \Prob\big(M - \E M \ge u\big) \le \exp\Big(-\frac{u^2}{2\cG\sigma^2}\Big), \qquad
  \Prob\big(M - \E M \le -u\big) \le \exp\Big(-\frac{u^2}{2\cG\sigma^2}\Big).
\]
(If $\sigma = 0$ then $\ip{x}{G} = 0$ a.s.\ for every $x\in\Xs$, $M = 0$ a.s., and both
probabilities vanish for $u > 0$.)
\end{theorem}

\begin{proof}
\uses{lem:compact_sup_subgaussian}
By Lemma~\ref{lem:compact_sup_subgaussian}, $M - \E M$ satisfies \texttt{HasSubgaussianMGF}
with constant $c := \cG\sigma^2$. The upper tail is the Chernoff bound
\texttt{HasSubgaussianMGF.measure\_ge\_le}: $\Prob(X \ge u) \le \exp(-u^2/(2c))$ for $u \ge 0$;
the lower tail is \texttt{HasSubgaussianMGF.measure\_le\_le} (LML; it is
\texttt{measure\_ge\_le} applied to $-X$, which is sub-Gaussian with the same constant by
\texttt{HasSubgaussianMGF.neg}). The degenerate case: $\Var\ip{x}{G} = x^\top Sx = 0$ forces
$\ip{x}{G} = 0$ a.s.\ for each $x$, hence for all $x$ in the countable dense set
simultaneously, hence $M = 0$ a.s.\ by Lemma~\ref{lem:sup_countable_dense}(ii).
\end{proof}

\begin{remark}[The classical statement]\label{rem:pc_borell_tis_classical}
The Borell--TIS inequality \cite[Theorem~2.1.1]{adler2007random} holds for every
almost surely bounded centered Gaussian process with $\cG$ replaced by $1$, and its proof
goes through the Gaussian isoperimetric inequality. Only the linear processes
$x\mapsto\ip{x}{G}$ over compact $\Xs \subset \R^d$ are needed in this blueprint, and only up
to the absolute constant $\cG$; the general statement is not pursued.
\end{remark}

\begin{corollary}[High-probability upper bound on the supremum]\label{cor:sup_bound_whp}
\uses{thm:borell_tis}
In the setting of Theorem~\ref{thm:borell_tis}, for every $\delta \in (0,1)$,
\[
  \Prob\Big(M \le \E M + \sigma\sqrt{2\cG\log(1/\delta)}\Big) \ge 1 - \delta, \qquad
  \Prob\Big(M \ge \E M - \sigma\sqrt{2\cG\log(1/\delta)}\Big) \ge 1 - \delta .
\]
\end{corollary}

\begin{proof}
\uses{thm:borell_tis}
Let $u := \sigma\sqrt{2\cG\log(1/\delta)} \ge 0$; then $u^2/(2\cG\sigma^2) = \log(1/\delta)$
and Theorem~\ref{thm:borell_tis} gives $\Prob(M - \E M \ge u) \le e^{-\log(1/\delta)} = \delta$.
The event $\{M \le \E M + u\}$ contains the complement of $\{M - \E M \ge u\}$ (in fact
equals the complement of $\{M - \E M > u\}$), so its probability is at least $1 - \delta$. The
second bound is the lower tail. If $\sigma = 0$ both events have probability one.
\end{proof}

\begin{corollary}[Symmetric index set]\label{cor:sup_bound_symmetric}
\uses{thm:borell_tis, lem:gauss_law_of_cov}
Let $\Xs \subset \R^d$ be nonempty and compact, $S \in \PSD$, $G \sim \Normal(0,S)$,
$\sigma^2 := \sup_{x\in\Xs}x^\top Sx$, and consider the difference set
$\Xs - \Xs := \{x - x' : x, x' \in \Xs\}$, which is nonempty and compact. Let
$D := \sup_{x,x'\in\Xs}\ip{x - x'}{G} = \sup_{z\in\Xs-\Xs}\ip{z}{G}$. Then
\begin{enumerate}
  \item[(i)] $\E D = 2\,\E\sup_{x\in\Xs}\ip{x}{G}$;
  \item[(ii)] $\sigma^2_{\Xs-\Xs} := \sup_{z\in\Xs-\Xs}z^\top Sz \le 4\sigma^2$;
  \item[(iii)] for every $\delta\in(0,1)$,
        $\Prob\big(D \le \E D + 2\sigma\sqrt{2\cG\log(1/\delta)}\big) \ge 1 - \delta$.
\end{enumerate}
\end{corollary}

\begin{proof}
\uses{cor:sup_bound_whp, lem:gauss_law_of_cov, lem:sup_countable_dense}
$\Xs - \Xs$ is the image of the compact $\Xs\times\Xs$ under the continuous map
$(x,x')\mapsto x - x'$, hence compact. (i) Pointwise,
$\sup_{x,x'}(\ip{x}{G} - \ip{x'}{G}) = \sup_x\ip{x}{G} + \sup_{x'}\ip{x'}{-G}$ (the two suprema
are over independent variables $x, x'$), and $-G \sim \Normal(0,S)$
(Lemma~\ref{lem:gauss_law_of_cov}(iv)), so $\E\sup_{x'}\ip{x'}{-G} = \E\sup_x\ip{x}{G}$; both
terms are integrable by Lemma~\ref{lem:sup_countable_dense}. (ii) For $z = x - x'$,
$z^\top Sz = \norm{S^{1/2}x - S^{1/2}x'}^2 \le (\norm{S^{1/2}x} + \norm{S^{1/2}x'})^2 \le (2\sigma)^2$.
(iii) Apply Corollary~\ref{cor:sup_bound_whp} to the index set $\Xs - \Xs$: with probability
at least $1 - \delta$, $D \le \E D + \sigma_{\Xs-\Xs}\sqrt{2\cG\log(1/\delta)} \le
\E D + 2\sigma\sqrt{2\cG\log(1/\delta)}$ by (ii). (If $\sigma_{\Xs - \Xs} = 0$ the bound holds
with probability one.)
\end{proof}

\textbf{Lean remark.} In the applications (Lemmas~\ref{lem:difference_process_concentration}
and~\ref{lem:region_phase1}) the index set is $\Xs - \Xs$ for the action set or for one region
of a partition, $S = \Sigma_T^{-1}$ for a fixed design, and $\sigma^2 = \max_{x}\norm{x}^2_{\Sigma_T^{-1}}$;
the compact set $\Xs - \Xs$ is \texttt{Set.image2 (· - ·) Xs Xs}, which is the image of the
product set $\Xs\times\Xs$ (\texttt{Set.prod}) under \texttt{fun p ↦ p.1 - p.2}
(\texttt{Set.image\_prod}); with \texttt{hK : IsCompact Xs} it is compact by
\texttt{(hK.prod hK).image continuous\_sub} (Mathlib \texttt{IsCompact.image}), or equivalently
\texttt{(hK.prod hK).image\_of\_continuousOn continuous\_sub.continuousOn}.
